\section{Introduction}
In this project, we study how the signatures of symmetry-protected topological (SPT) order generalize to a model protected by an exotic type of symmetry called a \textit{fusion category symmetry.} Let us first understand what fusion category symmetries are and how we can realize them in lattice models using generalized qudits.

Most of the symmetries which we study in quantum mechanics are described by a group $G$ and realized by a set of unitary or anti-unitary operators labeled by elements of the group $\{U_g:g\in G\}$. However, quantum mechanical systems can host much more general symmetries which are described not by groups, but by \textit{\textbf{fusion categories.}} 
\begin{definition}[Fusion category, informal]
    A fusion category $\mathcal{A}$ consists of a finite set of \textit{simple objects} $\{a_i\}$ and a product $\otimes:\mathcal{A}\times\mathcal{A}\rightarrow\mathcal{A}$ such that:
    \begin{enumerate}
        \item $a_i\times a_j=\sum_k N^{k}_{i,j} a_k$, $N_{i,j}^k\in\mathbb{N},$ \\and for each $a\in \mathcal{A}$:
        \item (Identity) There exists a trivial object $1\in\mathcal{A}$ such that $1\otimes a=a\otimes 1=a$,
        \item (Generalized Inverse) There exists an $\bar a\in\mathcal{A}$ such that $a\otimes\bar a=\bar a\otimes a=1+\cdots$.
    \end{enumerate}
\end{definition}
\textit{\textbf{Fusion category symmetries}} are symmetries of a quantum mechanical system realized by a set of operators labeled by the elements of a fusion category: $\{O_a:a\in\mathcal{A}\}$. Because of property 3, these symmetries are said to be \textit{non-invertible}, meaning there are some symmetry transformations which do not have a corresponding inverse in the sense in which we are familiar. To realize fusion category symmetry in a lattice model, we introduce \textit{group-valued qudits}.

\section{Group-Valued Qudits}
The generalized cluster state is a model defined on \textbf{\textit{group-valued qudits}}, which may be familiar as the degrees of freedom in Kitaev's quantum double model \cite{kitaev_fault-tolerant_2003-2}. To construct these group-valued qudits, we fix a finite group $G$ and work with the set of basis states
$\{\ket{g}:g\in G\}$
labeled by elements of the group. The model is constructed on a tensor product of this Hilbert space at each site. 

We can then define \textbf{\textit{generalized Pauli operators}} which act as group multiplication (Pauli $X$):
\begin{equation*}
    \lx_g|h\rangle=|gh\rangle,\quad \rx_g|h\rangle=|h\bar{g}\rangle,
\end{equation*}
and generalized phase (Pauli $Z$):
\begin{equation*}
\zt_{\Gamma}=
\begin{tikzpicture}[scale=2]
    \draw[red] (-0.5,0)--(0.5,0);
    \draw (0,-0.5) -- (0,0.5);
    \node[mpo] (t) at (0,0) {$\zt_\Gamma$};
  \end{tikzpicture}  
=\sum_{g\in G}\textcolor{red}{\Gamma(g)}\otimes \ketbra{g}{g}.
\end{equation*}
where $\Gamma$ is an irreducible representation (irrep) of $G$: $\Gamma:G\rightarrow\GL(\mathbb{C})$. Notice that $\zt_\Gamma$ is a matrix product operator (MPO). Indeed, tensor networks are the natural language to describe the generalized cluster state and its symmetries.

\section{Generalized Cluster State}
The generalized cluster state was introduced in \cite{brell_generalized_2015} and is given by
\begin{equation*}
    |\mathcal{C}\rangle=|G|^{-N/2}\sum_{\{g_i\}\in G}|g_1\rangle|g_1\bar{g_2}\rangle|g_2\rangle\cdots|g_{N-1}\bar{g_N}\rangle|g_N\rangle,
\end{equation*}
when placed on an open chain of $2N-1$ sites. It can also be expressed as the matrix product state 
\begin{equation*}
    \ket{\Cl}=\begin{tikzpicture}[scale=2]
    \draw (0,0)--(5.5,0);
    \foreach \x in {0.5,2,3.5,5}{
		     \draw (\x,0) -- (\x,0.5);
		    \node[odd] at (\x,0) {};
		  }
		  \foreach \x in {1.25,4.25}{
		     \draw (\x,0) -- (\x,0.5);
		    \node[even] at (\x,0) {};
		  }
		  \node[fill=white] at (2.75,0) {$\dots$};
		  \node[edge] at (0,0) {};
		  \node[edge] at (5.5,0) {};
  \end{tikzpicture} \ ,
\end{equation*}
where the MPS tensors are given by
\begin{equation*}
    \begin{tikzpicture}[scale=2]
    % \node[irrep] at (0.2,0) {};
    \draw (-0.5,0)--(0.5,0);
    \draw (0,0) -- (0,0.5);
    \node[odd] (t) at (0,0) {};
  \end{tikzpicture} 
\ =\sum_{g\in G}\ketbra{g}{g}\otimes\ket{g},\quad 
  \begin{tikzpicture}[scale=2]
    % \node[irrep] at (0.2,0) {};
    \draw (-0.5,0)--(0.5,0);
    \draw (0,0) -- (0,0.5);
    \node[even] (t) at (0,0) {};
  \end{tikzpicture} 
\ =\sum_{g\in G} \lx_g\otimes\ket{g}.
\end{equation*}
We can identify a family of generalized Pauli operators which stabilize $\ket{\Cl}$ and use these to construct a Hamiltonian which hosts $\ket{\Cl}$ as its ground state:
\begin{equation*}
    H_\Cl=-\frac{1}{|G|}\sum_{i\text{ odd}}\Bigg(\sum_{\Gamma\in\rep{G}}\Tr\left[\zt^{\dagger(i)}_\Gamma.\zt^{(i+1)}_\Gamma.\zt^{(i+2)}_\Gamma\right]d_\Gamma+\sum_{g\in G}\rx_g^{(i+1)}\lx_g^{(i+2)}\lx_g^{(i+3)}\Bigg).
\end{equation*}
We can identify two families of symmetries of this Hamiltonian:
\begin{equation*}
    \left\{\overleftarrow{A}_g = \prod_{i\text{ odd}}\rx_g^{(i)}:g\in G\right\},\quad \left\{ \hat B_\Gamma = \Tr\left[\prod_{i\text{ even}}\zt^{(i)}_\Gamma\right]:\Gamma\in\rep G\right\}.
\end{equation*}
These symmetries act on different sublattices and are therefore independent, so that the total symmetry is $G\times\rep G$. The generalized cluster state is also symmetric with respect to both of these symmetries. 
\vspace{1.5em}


\vspace{1em}


